\documentclass[10pt,a4paper]{amsart}
\usepackage[margin=19mm]{geometry}
\usepackage{amsmath,amssymb,mathtools}
\usepackage[scr=rsfs,cal=euler]{mathalfa}
\usepackage{xfrac}
\usepackage[unicode,hidelinks,bookmarks=false]{hyperref}
\newcommand{\ie}{i.e.\ }
\newcommand{\cf}{cf.\ }
\newcommand{\resp}{respectively\ }
\newcommand{\Subsection}{\subsection}
\providecommand{\nobreakdash}{\nobreak\hbox{-}}
\setlength{\emergencystretch}{2em}
\multlinegap0pt
\usepackage{amssymb}
\usepackage{xcolor}
\usepackage{mathtools}
\usepackage{enumerate}
\newtheorem{proposition}{Proposition}
\newcommand{\F}{\mathcal{F}}
\newcommand{\R}{\mathbb{R}}
\title{Sharp estimates for eigenvalues of localization operators before the plunge region}
\author{Aleksei Kulikov}
\date{Source: arXiv:2603.07407v1 (CC BY 4.0)}
\begin{document}
\maketitle

\noindent\emph{Source and licence.} Sharp estimates for eigenvalues of localization operators before the plunge region. Version arXiv:2603.07407v1 (CC BY 4.0), distributed by arXiv under the Creative Commons Attribution 4.0 International licence, http://creativecommons.org/licenses/by/4.0/, as recorded in the arXiv licence metadata for this exact version and shown on the abstract page https://arxiv.org/abs/2603.07407v1. Full licence notice: \texttt{licenses/arxiv-2603.07407-CC-BY-4.0.txt}.\par\smallskip

\noindent\textit{Editorial scope.} Counted unit: the article's proposition prop (source0.tex lines 469--475). The article's complete original proof of it is delivered with it (source0.tex lines 476--491, one \texttt{proof} environment of 16 lines). No mathematics is written, repaired or re-derived: the statement and the proof are the article's own lines, in the article's own notation, and no retained line is substituted or re-flowed.  No further source-local context is needed: every cross-reference of every retained line resolves inside the two delivered spans.  Nothing was omitted from any retained span, so the package carries no omission record.  No retained line contains a citation, so the wrapper carries no bibliography and no bibliography entry.  No retained line contains a figure environment, an image-inclusion command or drawing code, so no image is delivered and the package declares no asset dependency.  Editorial formatting only, in addition: an AMS \texttt{amsart} preamble that restates the article's own declarations and macros used by the retained lines, namely the environments \texttt{proposition} on separate counters, together with the standard packages those lines need.  The retained environments print in this document as Proposition 1 in order of appearance, whereas the article prints the same items with the numbering of its own full document.  The complete native source of the article as distributed in the arXiv e-print for this exact version is bundled unchanged as \texttt{sources/195941/source0.tex}.
\begin{proposition}\label{function prop}
Let $T, W>0$ be such that $TW \ge 1$. The function $$f_{T, W} =\F \left(e^{-\pi x^2 \frac{W}{T}}\chi_{[-\frac{T}{2},\frac{T}{2}]}(x)\right)$$ satisfies
$$\int_{|t| >\frac{W}{2}} |f_{T,W}(t)|^2dt \le \sqrt{\frac{T}{W}}10e^{-\frac{\pi TW}{2}},$$
$$\int_{\R} |f_{T,W}(t)|^2dt\le \sqrt{\frac{T}{W}},$$
$$f_{T, W}(0)\ge \frac{1}{10}\sqrt{\frac{T}{W}},$$
$$|f_{T,W}(t)|\le f_{T,W}(0), t\in\R.$$
\end{proposition}

\begin{proof}
First of all we notice that all of the estimates respect the linear change of variables $t\to qt$, so we can without loss of generality assume that $T = W\ge 1$. In this case, since $e^{-\pi x^2}$ is its own Fourier transform, we get
 $$f_{T,T} = \F \left(e^{-\pi x^2} \chi_{[-\frac{T}{2},\frac{T}{2}]}\right) = e^{-\pi x^2} - \F \left(e^{-\pi x^2} \chi_{\R\backslash [-\frac{T}{2},\frac{T}{2}]}\right).$$
This means, using the triangle inequality and the Plancherel theorem, that 
$$\int_{|t| > \frac{W}{2}} |f_{T, T}(t)|^2 dt \le 8\int_{\frac{T}{2}}^\infty e^{-2\pi x^2}dx.$$
This is the tail of the Gaussian for which very precise estimates are known, but for us a simple bound is sufficient:
$$\int_{\frac{T}{2}}^\infty e^{-2\pi x^2}dx \le \frac{2}{T}\int_{\frac{T}{2}}^\infty xe^{-2 \pi x^2}dx = \frac{1}{2\pi T}e^{-\frac{\pi T^2}{2}} \le e^{-\frac{\pi T^2}{2}},$$
where in the last step we used $T \ge 1$.

The second estimate is the easiest, as it follows directly from the Plancherel theorem:
$$\int_\R |f_{T,T}(t)|^2dt = \int_{-\frac{T}{2}}^{\frac{T}{2}} e^{-2\pi x^2}dx \le \int_\R e^{-2\pi x^2}dx =\frac{1}{\sqrt{2}} < 1.$$

For the value at $0$ we have
$$f_{T,T}(0) = \int_{-\frac{T}{2}}^{\frac{T}{2}}e^{-\pi x^2}dx \ge \int_{-\frac{1}{2}}^{\frac{1}{2}}e^{-\pi x^2}dx \ge e^{-\frac{\pi}{4}} \ge \frac{1}{e} > \frac{1}{10}.$$
Finally, the estimate $|f_{T,T}(t)|\le f_{T,T}(0)$ follows from the fact that the Fourier transform of $f_{T,T}$ is non-negative.
\end{proof}

\end{document}
